\documentclass[11pt]{article}

% Encoding and font setup (improved for pdftotext extraction quality).
\usepackage[a4paper,margin=2.5cm]{geometry}
\usepackage[utf8]{inputenc}
\usepackage[T1]{fontenc}
\usepackage{lmodern}
\usepackage{textcomp}
\input{glyphtounicode}
\pdfgentounicode=1

\usepackage{amsmath,amssymb,amsthm,mathtools}
\usepackage{authblk}
\usepackage[final,babel=true,activate={true,nocompatibility},protrusion=true,expansion=false]{microtype}
\usepackage{booktabs}
\usepackage{longtable}
\usepackage{array}
\usepackage{enumitem}
\usepackage{xcolor}
\usepackage{hyperref}

\hypersetup{
  colorlinks=true,
  linkcolor=blue,
  citecolor=blue,
  urlcolor=blue,
  pdftitle={A conjectural cardinality envelope for the condensed classifying anima of spectral infinity-topoi},
  pdfauthor={Vera Gomez, Francisco Javier},
  pdfsubject={Condensed mathematics, model theory, set theory},
  pdfkeywords={condensed mathematics, condensed homotopy type, Lascar group, infinity-topos, cardinal arithmetic, LLM-assisted research}
}

% Theorems
\theoremstyle{definition}
\newtheorem{conjecture}{Conjecture}
\newtheorem*{question}{Question}
\newtheorem*{convention}{Convention}
\theoremstyle{plain}
\newtheorem{theorem}{Theorem}
\newtheorem{proposition}[theorem]{Proposition}
\newtheorem{lemma}[theorem]{Lemma}
\newtheorem{corollary}[theorem]{Corollary}
\theoremstyle{remark}
\newtheorem*{remark}{Remark}

% Convenience
\newcommand{\Cond}{\mathrm{Cond}}
\newcommand{\Pt}{\mathrm{Pt}}
\newcommand{\Bcat}{\mathrm{B}}
\newcommand{\Sh}{\mathrm{Sh}}
\newcommand{\Map}{\mathrm{Map}}
\newcommand{\GalL}{\mathrm{Gal}_{L}}
\newcommand{\Galabs}{\mathrm{Gal}}
\newcommand{\Frhat}{\widehat{\mathrm{Fr}}}
\newcommand{\Zhat}{\widehat{\mathbf{Z}}}
\newcommand{\Aone}{\mathbf{A}^{1}}
\newcommand{\Pone}{\mathbf{P}^{1}}
\newcommand{\C}{\mathbf{C}}
\newcommand{\Q}{\mathbf{Q}}
\newcommand{\R}{\mathbf{R}}
\newcommand{\continuum}{\mathfrak{c}}

\setlist[itemize]{nosep,leftmargin=*}
\setlist[enumerate]{nosep,leftmargin=*}

\title{A conjectural cardinality envelope for the\\
  condensed classifying anima of spectral $\infty$-topoi:\\
  a synthesis across geometric and analogical instances}

\author[1]{Francisco Javier Vera G\'omez}
\affil[1]{Independent researcher\\
\href{https://orcid.org/0009-0001-3516-5871}{ORCID: 0009-0001-3516-5871}\\
Repository:
\href{https://gitlab.loneorc.com/research/axiom-explorer}{gitlab.loneorc.com/research/axiom-explorer}\\
Version: v1.2-preprint}

\date{\today}

\begin{document}
\maketitle

\begin{abstract}
We formulate a candidate cardinality envelope for the underlying group at section
level of the fundamental group of the condensed classifying anima of a spectral
$\infty$-topos: with $\kappa_{\mathcal X}$ a size invariant defined in
\S\ref{sec:size}, we conjecture
$|\,\pi_1(\Bcat\,\Pt^{\mathrm{coh}}(\mathcal X))(\ast)\,| \le 2^{\kappa_{\mathcal X}}$.
We test this against two direct geometric attestations (\cite{HHLMMW2025}, on
$\Pone_{\C}$ and $\Pone_{\Q}$) and discuss related cardinal phenomena in
model theory (Lascar groups via the unifying classifying-anima construction
recently announced by \cite{Haine2026classifying}) and set theory (Whitehead's
problem in condensed mathematics, after Clausen-Scholze and follow-ups).
The unifying classifying-anima construction is itself an announced result; the
envelope as a uniform statement \emph{across} the unification is, to the best
of our knowledge, not formulated as such in any single source. We present the
candidate as a falsifiable conjecture, with explicit confidence levels per
claim, an honest declaration of author position and production method, and an
explicit invitation to specialists to confirm, refute, or locate it in
existing literature. \emph{This preprint is the output of an LLM-assisted
synthesis workflow, axiom-explorer; see \S\ref{sec:method} for full
disclosure.}
\end{abstract}

\section{Introduction and statement}\label{sec:intro}

Let $\mathcal{X}$ be a coherent $\infty$-topos. Following Lurie's work on
ultracategories, one associates to $\mathcal{X}$ a \emph{condensed
$\infty$-category of points} $\Pt(\mathcal{X})$, defined on extremally
disconnected profinite sets $K$ by
$\Pt(\mathcal{X})(K) := \mathrm{Fun}^{\ast}(\mathcal{X}, \Sh(K))$, the
$\infty$-category of left-exact left adjoints
$f^{\ast}: \mathcal{X} \to \Sh(K)$. Restricting to $f^{\ast}$ that
\emph{preserve coherent objects} gives the condensed $\infty$-category of
\emph{coherent points} $\Pt^{\mathrm{coh}}(\mathcal{X})$ of
\cite{BarwickGlasmanHaine2018}. We write
$\Bcat\,\Pt^{\mathrm{coh}}(\mathcal{X})$ for the condensed classifying anima of
this condensed $\infty$-category.

\subsection{Size conventions}\label{sec:size}

Earlier formulations of this work used the cardinality of the underlying set
of objects of $\Pt(\mathcal X)(\ast)$ as a complexity parameter. That is not
invariant under equivalence of $\infty$-categories. Following a reviewer's
correction, we use:

For a topological group $G$, let $\mathrm{wt}(G)$ denote the topological
weight of $G$ (the minimum cardinality of a basis of open neighbourhoods of
the identity). For a profinite group $G$, $\mathrm{wt}(G)$ equals the minimum
cardinality of a generating family of $G$ as a topological group, and one has
$|G| \le 2^{\mathrm{wt}(G)}$ in general, with equality for free profinite
groups.

\begin{convention}\label{conv:kappa}
For a spectral $\infty$-topos $\mathcal X$, let
\begin{equation}\label{eq:kappa}
\kappa_{\mathcal X} \;:=\; \max\!\left(
  \aleph_0,\;
  |\pi_0\,\Pt(\mathcal X)(\ast)|,\;
  \sup_{x \in \pi_0\,\Pt(\mathcal X)(\ast)} \mathrm{wt}\bigl(\pi_1(\Pt(\mathcal X)(\ast),\,x)\bigr)
\right).
\end{equation}
We use $\continuum := 2^{\aleph_0}$ throughout.
\end{convention}

The third term measures the \emph{topological weight} of the loop groups at
points (typically profinite generating size), not the cardinality of the
underlying group. This distinction is crucial: a free profinite group on
a set $S$ has weight $|S|$ but cardinality up to $2^{|S|}$. Using cardinality
in \eqref{eq:kappa} would force the bound to grow one full beth level
beyond what the geometric attestations actually realise; using weight
recovers the regime in which the conjecture is meant to hold.

The conjecture below is invariant under equivalence of $\infty$-categories
of points since each ingredient on the right of \eqref{eq:kappa} is.

\begin{remark}\label{rem:kappa-coarse}
Convention \ref{conv:kappa} is one specific choice. A sharper invariant
might further weight the higher Postnikov data of $\Pt(\mathcal X)$. Refining
the precise size invariant is left open as part of the conjecture's
calibration.
\end{remark}

\begin{remark}[Reviewer correction, v1.1 \textrightarrow{} v1.2]
An earlier draft used $\sup |\pi_0\Map(x,y)|$ as the third term in
\eqref{eq:kappa}. A reviewer pointed out that for the geometric \'etale
$\infty$-topos this term equals (up to $\pi_0$) the cardinality of the
absolute Galois group of a generic point, which can already be as large as
$2^{|k|}$, breaking the geometric attestations of \S\ref{sec:direct}. The
weight-based formulation above is the corresponding fix.
\end{remark}

\subsection{Statement}

\begin{conjecture}[Cardinality envelope, upper bound]\label{conj:envelope}
For $\mathcal{X}$ a spectral $\infty$-topos and $\kappa_{\mathcal X}$ as in
Convention \ref{conv:kappa},
\begin{equation}
\bigl| \pi_1(\Bcat\,\Pt^{\mathrm{coh}}(\mathcal{X}))(\ast) \bigr|
\;\le\; 2^{\kappa_{\mathcal X}}.
\end{equation}
\end{conjecture}

The conjecture asserts a one-line cardinal bound, no more. Its
\emph{attainment} is treated separately:

\begin{question}[Saturation]\label{q:saturation}
Under what hypotheses on $\mathcal X$ is the upper bound of
Conjecture \ref{conj:envelope} attained? More specifically: for which
$\mathcal X$ does $\Pt(\mathcal X)(\ast)$ admit a $\kappa_{\mathcal X}$-large
family of distinct localisation directions whose closure under intersection
again has cardinality $\kappa_{\mathcal X}$, and does this combinatorial
condition force saturation?
\end{question}

The earlier draft of this work conflated the upper bound with a saturation
claim. The reviewer pointed out (correctly) that the saturation claim is much
stronger and much less well supported. We separate them here.

\subsection{Why this conjecture}

The conjecture sits on top of a structural unification announced by Haine
\cite[Theorem 0.3]{Haine2026classifying}: for spectral $\infty$-topoi
$\mathcal{X}$, the inclusion of condensed $\infty$-categories
$\Pt^{\mathrm{coh}}(\mathcal{X}) \hookrightarrow \Pt(\mathcal{X})$ induces an
equivalence on condensed classifying anima. Under this equivalence, two
superficially unrelated invariants become special cases of the same
construction:

\begin{itemize}
\item The pro\'etale fundamental group of a scheme of Bhatt-Scholze
\cite{BhattScholze2015}, via the \'etale $\infty$-topos
$\mathcal{X} = X_{\mathrm{\acute{e}t}}$ --- refined to the condensed
fundamental group $\pi_1^{\mathrm{cond}}(X)$ of \cite{HHLMMW2025}.
\item The Lascar group $\GalL(T)$ of a complete first-order theory $T$, via
the classifying $\infty$-topos of $T$ --- recovered as $\pi_1$ of the
condensed classifying anima of $\mathrm{Mod}_T$ in forthcoming work of Haine,
Damaj and Zhang announced in \cite[\S 0]{Haine2026classifying}. The classical
unconditional identification $\GalL(T) \simeq \pi_1(|\mathrm{Mod}_T|)$ is due
to Campion, Cousins and Ye \cite{CampionCousinsYe2021}; the condensed-group
upgrade is the announced extension.
\end{itemize}

The conjecture asks whether a uniform cardinality bound holds across this
unification.

\section{Confidence ladder}\label{sec:confidence}

We mark each substantive claim with a confidence level so the reader can
adjust trust accordingly:

\begin{itemize}
\item \textbf{L0 --- Verified mechanically}: directly checked in source code,
  formal proof, or finite enumeration.
\item \textbf{L1 --- Strong}: a standard, named result in the cited literature.
\item \textbf{L2 --- Plausible}: argued or supported by the cited evidence; not
  directly verified by us.
\item \textbf{L3 --- Speculative}: a guess emerging from the synthesis, listed
  as open question.
\end{itemize}

Conjecture \ref{conj:envelope} carries confidence \textbf{L2} for the upper
bound. Question \ref{q:saturation} (saturation) is \textbf{L3}.

\section{Direct evidence: geometric attestations}\label{sec:direct}

For $X$ a quasi-compact quasi-separated scheme, the condensed homotopy type
$\Pi_\infty^{\mathrm{cond}}(X)$ is the condensed classifying anima of
$\Pt^{\mathrm{coh}}(X_{\mathrm{\acute{e}t}})$ \cite[\S 3]{HHLMMW2025}.

The relevant attestations of Conjecture \ref{conj:envelope} are explicit in
\cite[Remark 7.14, Example 7.9]{HHLMMW2025}.

\begin{theorem}[Geometric attestations, after \cite{HHLMMW2025}]\label{thm:geomatt}
Working over the complex numbers and over the rationals respectively,
\begin{equation}
\bigl| \pi_1^{\mathrm{cond}}(\Pone_{\C})(\ast) \bigr| = 2^{\continuum},
\qquad
\bigl| \pi_1^{\mathrm{cond}}(\Pone_{\Q})(\ast) \bigr| \le \continuum.
\end{equation}
\end{theorem}

\paragraph{Computing $\kappa$ for $\Pone_{\C}$.}
For $\Pone_{\C}$, the closed points are indexed by $\C \sqcup \{\infty\}$, so
$|\pi_0\Pt(\Pone_{\C})(\ast)| = \continuum$. The loop group at the generic
point is the absolute Galois group of $\C(T)$, a free profinite group on a
set of size $\continuum$ \cite[Theorem 2]{Douady1964}, of \emph{topological
weight} $\continuum$ even though its underlying cardinality is
$2^{\continuum}$. Hence $\kappa_{\Pone_{\C}} = \continuum$ under
Convention \ref{conv:kappa}, the envelope $2^{\kappa}$ is $2^{\continuum}$,
and the attested $|\pi_1^{\mathrm{cond}}(\Pone_{\C})(\ast)| = 2^{\continuum}$
\textbf{attains} it.

\paragraph{Computing $\kappa$ for $\Pone_{\Q}$.}
For $\Pone_{\Q}$, closed points form a countable set and the absolute
Galois group of $\Q(T)$ is a free profinite group on a countable index set,
of weight at most $\aleph_0$. Hence $\kappa_{\Pone_{\Q}} = \aleph_0$, the
envelope $2^{\kappa}$ is $\continuum$, and the attested upper bound
$|\pi_1^{\mathrm{cond}}(\Pone_{\Q})(\ast)| \le \continuum$ sits inside it.
Whether the bound is also \emph{attained} for $\Pone_{\Q}$ is, as far as we
know, open: it bears on Question \ref{q:saturation}.

\subsection{The mechanism in \cite{HHLMMW2025}}

The proof of \cite[Proposition 7.2]{HHLMMW2025} proceeds by a short exact
sequence
\begin{equation}\label{eq:ses}
1 \;\to\; N_S \;\to\; \Frhat_{\C} \;\to\;
\pi_1^{\mathrm{cond}}(\Aone_{\C} \setminus S, \overline{\eta})(\ast)
\;\to\; 1,
\end{equation}
where $\Frhat_{\C}$ is the free profinite group on $\C$ and $N_S$ is an
explicit abstract normal subgroup. Two inputs control the cardinality of the
quotient:

\begin{itemize}
\item The cardinality of the free profinite group:
\begin{equation}\label{eq:Frbound}
|\Frhat_{S}| \;\le\; 2^{\max(|S|,\,\aleph_0)}
\end{equation}
for any non-empty set $S$ (essentially because $\Frhat_S$ is the inverse limit
of finite quotients indexed by finite subsets of $S$, and the cardinality is
controlled by the number of continuous maps to finite groups).
\item A natural surjection onto the quotient
$\prod_{a \in \C} \Zhat \,/\, \bigoplus_{a \in \C} \Zhat$, which has
cardinality $2^{\continuum}$.
\end{itemize}

For $k = \C$: \eqref{eq:Frbound} gives $|\Frhat_{\C}| \le 2^{\continuum}$, the
quotient surjects onto a group of size $2^{\continuum}$, and the two bounds
pinch the answer to $2^{\continuum}$. For $k = \Q$: \eqref{eq:Frbound} gives
$|\Frhat_{\Q}| \le 2^{\aleph_0} = \continuum$, the analogous lower bound is at
most $\continuum$, and the upper bound of $\continuum$ stands.

\begin{remark}
The earlier draft of this preprint stated $|\Frhat_S| \le |S|^{\aleph_0}$,
which is incorrect for $|S|$ uncountable: it would give $|\Frhat_{\C}| \le
\continuum$, in conflict with the attested
$|\pi_1^{\mathrm{cond}}(\Pone_{\C})| = 2^{\continuum}$. The correct bound
\eqref{eq:Frbound} resolves the inconsistency. We thank a reviewer for flagging
this.
\end{remark}

\section{Related cardinal phenomena}\label{sec:analogical}

Beyond the two direct attestations, several lines of recent work concern
cardinal regimes that are \emph{analogous} to the envelope of
Conjecture \ref{conj:envelope}, in the sense that the same cardinal
$2^{\kappa}$-shape appears for related (but not identical) invariants. We
present these as motivating context, \emph{not} as direct attestations of the
conjecture.

\subsection{Model theory: the Lascar group}\label{sec:modeltheory}

For a complete first-order theory $T$, the Lascar group
$\GalL(T) = \mathrm{Aut}(\mathfrak C / \mathrm{acl}^{\mathrm{eq}}(\emptyset))\,/\,\mathrm{Gal}_{L,0}(T)$
is a quasi-compact topological group attached to a sufficiently saturated
``monster'' model $\mathfrak C$ \cite[Ch.~6]{Casanovas2011}. It is independent
of the choice of $\mathfrak C$ above some saturation cardinal.

The classical model-theoretic upper bound is the sharp inequality
\begin{equation}\label{eq:lascar-bound}
|\GalL(T)| \;\le\; 2^{|T|},
\end{equation}
where $|T|$ is the cardinality of the language plus parameters
\cite[Theorem 1.4]{CampionCousinsYe2021}. For countable $T$ this gives
$|\GalL(T)| \le \continuum$.

Campion, Cousins, and Ye \cite{CampionCousinsYe2021} establish the
unconditional identification
\begin{equation}
\GalL(T) \;\cong\; \pi_1\bigl(|\mathrm{Mod}_T|\bigr),
\end{equation}
the fundamental group of the geometric realisation of the category of
models. The condensed-group upgrade of this identification, asserted in
\cite[\S 0]{Haine2026classifying} as joint forthcoming work of Haine, Damaj
and Zhang, would place $\GalL(T)$ inside the framework of
Conjecture \ref{conj:envelope}.

What is true \emph{classically} about $|\GalL(T)|$ in two specific cases:

\begin{itemize}
\item For ACF$_0$ (algebraically closed fields of characteristic $0$), in
the formulation without naming constants for $\overline{\Q}$,
$\GalL(\mathrm{ACF}_0) \cong \Galabs(\overline{\Q}/\Q)$
\cite[Example 4.6]{CampionCousinsYe2021}, a profinite group of cardinality
$\continuum$. With the natural model-theoretic size invariant
$\kappa = |T| = \aleph_0$, the envelope $2^{\kappa} = \continuum$ is
\textbf{attained} cardinally.
\item For DLO (dense linear orders without endpoints),
$|\mathrm{Mod}(\mathrm{DLO})|$ is contractible
\cite[Example 4.4]{CampionCousinsYe2021}, hence $\GalL(\mathrm{DLO})$ is
trivial. This sits trivially inside the envelope. (DLO is the prototype of
an \emph{unstable} theory \cite{nLab-DLO}; we retract the v1 claim that
DLO is ``simple''.)
\end{itemize}

What we do \emph{not} attempt to claim, after the reviewer's correction:

\begin{itemize}
\item Specific cardinality bounds on $\GalL(\mathrm{Th}(\Q_p))$ or
on Lascar groups of ``wild'' theories at saturation. These cases plausibly
illustrate the envelope but require precise references to known computations
which we have not located. We do not include them in any tabulated evidence.
\end{itemize}

The classical bound \eqref{eq:lascar-bound} is tight up to a factor of $\le c$
in many natural cases (\cite[\S 4]{CampionCousinsYe2021}). The proposed
envelope of Conjecture~\ref{conj:envelope} is not meant to improve this
classical model-theoretic bound in isolation; its point is to ask whether the
same size principle persists \emph{uniformly across} the geometric and
model-theoretic sides of the announced Haine-Damaj-Zhang classifying-anima
identification.

\subsection{Set theory: condensed mathematics and forcing}

A second line of recent work concerns condensed mathematics and set theory.
Clausen and Scholze \cite{ClausenScholze2019Lectures} resolved Whitehead's
problem affirmatively in $\Cond(\mathrm{Ab})$ for discrete condensed abelian
groups, in contrast with Shelah's classical independence result for ordinary
abelian groups. Two recent papers give alternative proofs:

\begin{itemize}
\item Bergfalk, Lambie-Hanson, and \v Saroch \cite{BLHS2023} give a
forcing-theoretic argument, interpreting condensed sets as ``generically
invariant information''.
\item Bannister and Basak \cite{BannisterBasak2026} exhibit a geometric
morphism from the Grothendieck topos representing the Solovay model to the
$\kappa$-pyknotic sets of \cite{BarwickGlasmanHaine2018} (and Clausen-Scholze).
\end{itemize}

These results govern the \emph{derived} structure of $\Cond(\mathrm{Ab})$ and
its sensitivity to forcing axioms; they are \emph{not} direct computations of
$|\pi_1(\Bcat\,\Pt^{\mathrm{coh}}(\mathcal X))(\ast)|$. The cardinal regime
that emerges from these papers --- objects controlled by $|\R| = \continuum$
and at most $2^{\continuum}$ --- is consistent with the envelope of
Conjecture \ref{conj:envelope}, but as analogical evidence rather than direct
instantiation.

The fine-grained set-theoretic constraints of Bergfalk and Lambie-Hanson
\cite{BergfalkLH2024} on derived limits $\lim^n A_{\kappa,\lambda}$ are
strictly stronger than --- and not implied by --- a cardinality envelope on
the underlying group. We discuss this in \S\ref{sec:scope}.

\section{Summary table}

\begin{center}
\begin{longtable}{@{}p{2.4cm} p{4.6cm} p{2.0cm} p{3.4cm} p{2.0cm}@{}}
\toprule
\textbf{Type} & \textbf{Instance} & \textbf{$\kappa_{\mathcal X}$} & \textbf{Bound} & \textbf{Source} \\
\midrule
Direct (\S\ref{sec:direct}) & $\pi_1^{\mathrm{cond}}(\Pone_{\C})(\ast)$ & $\continuum$ & $= 2^{\continuum}$ (attained) & \cite[Rmk 7.14]{HHLMMW2025} \\
Direct (\S\ref{sec:direct}) & $\pi_1^{\mathrm{cond}}(\Pone_{\Q})(\ast)$ & $\aleph_0$ & $\le \continuum$ & \cite[Rmk 7.14]{HHLMMW2025} \\
\addlinespace
Analogical (\S\ref{sec:modeltheory}) & $\GalL(\mathrm{ACF}_0) \simeq \Galabs(\overline{\Q}/\Q)$ & $\aleph_0$ & $= \continuum$ (\textbf{attained}) & \cite[Ex.~4.6]{CampionCousinsYe2021} \\
Analogical (\S\ref{sec:modeltheory}) & $\GalL(\mathrm{DLO})$ & $\aleph_0$ & trivial (within envelope) & \cite[Ex.~4.4]{CampionCousinsYe2021} \\
\addlinespace
Analogical (set-th.) & Whitehead in $\Cond(\mathrm{Ab})$ for discrete & --- & cardinal regime $\le 2^{\continuum}$ & \cite{ClausenScholze2019Lectures}, \cite{BLHS2023} \\
Analogical (set-th.) & Solovay $\to$ $\kappa$-pyknotic morphism & --- & cardinal regime $\le 2^{\continuum}$ & \cite{BannisterBasak2026} \\
\bottomrule
\end{longtable}
\end{center}

The two \emph{direct} entries test the upper bound of
Conjecture \ref{conj:envelope}; the four \emph{analogical} entries describe
related cardinal regimes but do not directly instantiate the conjecture, by
the reviewer's correct observation. We do not claim ``seven attestations'',
contrary to an earlier draft.

\section{Structural argument (proof sketch)}\label{sec:sketch}

We sketch why the envelope is plausible structurally. \emph{This is not a
proof.}

Let $\mathcal{X}$ be spectral with $\kappa = \kappa_{\mathcal X}$ as in
\eqref{eq:kappa}. Two structural inputs control the size of
$\pi_1(\Bcat\,\Pt^{\mathrm{coh}}(\mathcal{X}))(\ast)$:

\begin{enumerate}
\item \textbf{An ambient profinite group}, in the form of an absolute Galois
or automorphism group attached to the data of $\Pt(\mathcal X)(\ast)$, has
\emph{topological weight} at most $\kappa$ (controlled by the third term of
\eqref{eq:kappa}), hence underlying cardinality at most $2^{\kappa}$ by
\eqref{eq:Frbound}. The argument runs through the
weight/generation-size invariant, not through the cardinality of the
group of morphisms in $\Pt(\mathcal X)$.
\item \textbf{Quotienting by an abstract normal closure} can only shrink the
group. Hence $|\pi_1| \le 2^{\kappa}$.
\end{enumerate}

The lower bound side --- the existence of a natural product family indexed by
$\kappa$ many points whose quotient by the direct sum reaches $2^{\kappa}$ ---
is the point flagged by the reviewer as the \emph{most fragile heuristic} of
the generalisation. For $X = \Pone_{\C}$ over the algebraically closed field $\C$,
this lower bound input is realised by the construction of
\cite[\S 7.1]{HHLMMW2025}, which is what makes $\Pone_{\C}$ saturate.
The analogous question for $\Pone_{\Q}$ remains open as far as we know. Whether
the analogous lower-bound input is available for general spectral
$\mathcal X$ is precisely the content of Question \ref{q:saturation}, and we
\emph{do not} claim it.

\begin{remark}
The structural argument as presented gives only the upper bound of
Conjecture \ref{conj:envelope}. The saturation claim of
Question \ref{q:saturation} requires a much sharper input that we have not
located in the literature in the level of generality we would need.
\end{remark}

\section{Scope of the candidate}\label{sec:scope}

We deliberately limit the scope of Conjecture \ref{conj:envelope} to a
\textbf{cardinality} statement on the \emph{underlying group at section
level}. The conjecture does \textbf{not} claim:

\begin{itemize}
\item \textbf{Bounded derived dimensions.} Bergfalk-Lambie-Hanson
\cite[Theorem D]{BergfalkLH2024} show that for $S, T$ extremally disconnected
\v{C}ech-Stone compactifications of discrete sets of cardinality $\ge
\aleph_\omega$, the constant sheaf on $S \times T$ has \emph{infinite}
injective dimension; in consequence, products of compact projective condensed
anima are not compact. The cardinality envelope is not in tension with this
--- it bounds underlying cardinalities, not Ext-tower heights.
\item \textbf{ZFC unconditionality} on the additivity of strong homology or
on the vanishing of $\lim^n A_{\kappa,\lambda}$ for $\kappa = \lambda =
\aleph_0$. These properties are genuinely sensitive to set-theoretic axioms
\cite[Theorem 2.10]{BergfalkLH2024}.
\end{itemize}

The cardinality envelope is, in this sense, a \emph{first-order} structural
observation --- useful as an upper bound, but not the deepest layer of
structure available in condensed mathematics.

\section{Falsifiability}\label{sec:falsifiability}

Conjecture \ref{conj:envelope} is sharply falsifiable. A falsifier is an
attested computation of
$|\pi_1(\Bcat\,\Pt^{\mathrm{coh}}(\mathcal X))(\ast)|$ exceeding
$2^{\kappa_{\mathcal X}}$. Concretely:

\begin{itemize}
\item A smooth qcqs scheme $X$ over an infinite field $k$ with
$|\pi_1^{\mathrm{cond}}(X_k)(\ast)|$ exceeding $2^{\kappa_{X_{\mathrm{\acute{e}t}}}}$,
with the right-hand side computed using \eqref{eq:kappa} on the \'etale
$\infty$-topos of $X$.
\item Conditional on the announced Haine-Damaj-Zhang condensed-group upgrade
of \cite{CampionCousinsYe2021}: a complete theory $T$ for which the condensed
$\GalL(T)$ exceeds $2^{\kappa_{T}}$ for the appropriate $\kappa_T$ derived
from the classifying $\infty$-topos.
\end{itemize}

We note that the falsifier of \emph{Question \ref{q:saturation}} is
different: a non-attaining $\mathcal X$ with the combinatorial structure
expected to force saturation. Falsifying saturation is much easier than
falsifying the upper bound, and would refine, not refute, the upper-bound
conjecture.

\begin{remark}
The earlier draft of this preprint listed
``$|\GalL(T)|$ exceeding $2^{|\mathrm{monster}|}$'' as a falsifier. The
reviewer pointed out that $\GalL(T)$ is always a quotient of a subgroup of
$\mathrm{Aut}(\mathfrak C)$, so this bound is essentially tautological and
not a useful falsifiability criterion. We retract that formulation; the
correct falsifier uses $\kappa_T$ from \eqref{eq:kappa}, not the cardinality
of the monster.
\end{remark}

\section{Author position and production method}\label{sec:method}

\subsection{Author position}

This preprint is produced by a non-specialist in any of the three branches
it engages: the author works at the boundary of formal methods and software
engineering. The mathematical content should be read as a \textbf{synthesis}
of recently-published results, not as expert exposition. Domain experts in
condensed mathematics, model theory, or set theory are explicitly invited to
flag misreadings of any cited literature.

\subsection{Production method}

The paper is the output of an autonomous experimental harness called
\emph{axiom-explorer} \cite{AxiomExplorerRepo}. Its workflow:

\begin{itemize}
\item Selects a small set of ``modern productive'' axiomatic seeds across
distinct branches.
\item Performs a systematic bibliometric cross-search over arXiv for each
pair of seeds (six binary combinations from four seeds; in this run,
Univalence Axiom $\times$ Condensed Mathematics $\times$ Perfectoid Spaces
$\times$ Synthetic Ricci).
\item Inventories the formal state of each seed in Mathlib4 and adjacent
Lean/Agda projects.
\item For each surviving cross-pair, builds a dossier integrating literature,
finite-case sanity checks (Z3, SymPy), and where applicable, a Lean 4
axiomatic skeleton verifying joint type consistency of the axioms involved.
\item Produces final per-phase reports and a synthesis document with
explicit confidence levels.
\end{itemize}

The harness is LLM-assisted (the orchestrator is an LLM agent operating
under explicit stop rules, and the expert-curator role of filtering
relevance is played jointly by the LLM and the human author). All inputs and
outputs of the run are versioned at the repository in
\cite{AxiomExplorerRepo} under a CC-BY-4.0 license.

The candidate of this preprint emerged from the harness's Phase 3 deep dive
on the Univalence $\times$ Condensed and Condensed $\times$ Perfectoid pairs,
followed by manual cross-reading that surfaced \cite{Haine2026classifying},
\cite{BergfalkLH2024}, \cite{BLHS2023} and \cite{BannisterBasak2026}. The
\textbf{unification} of the geometric and model-theoretic sides is announced
in those papers; it is not original to this work. The \textbf{uniform
cardinality envelope} \emph{across} the unification is, to the best of our
search, not formulated as such in any single source --- and is the synthesis
contribution of this preprint.

\section{Acknowledgments}

We thank the authors of all cited works for the source material. The Lean 4
axiomatic skeleton of Synthetic Stone Duality \cite{CCGM2024} in the
accompanying repository was built in dialogue with that paper's stated
axioms.

We are grateful to multiple anonymous AI reviewers whose commentary on
successive drafts led to the substantive corrections embodied in v1.1 and
v1.2.

The v1.1 corrections, prompted by a first round of review, included:
separating the upper bound from the saturation claim (\S\ref{sec:intro});
introducing an equivalence-invariant size invariant $\kappa_{\mathcal X}$
(an earlier first attempt, in v1.1, was further refined in v1.2); fixing
the cardinality bound on free profinite groups \eqref{eq:Frbound};
correcting the model-theoretic discussion of $\GalL(\mathrm{ACF}_0)$ and
DLO (\S\ref{sec:modeltheory}); demoting the set-theoretic discussion from
``attestation'' to ``related cardinal phenomena'' (\S\ref{sec:analogical});
reformulating the falsifiability section to avoid a near-tautological bound
(\S\ref{sec:falsifiability}); and toning the title from \emph{cardinality
envelope} to \emph{conjectural cardinality envelope}.

The v1.2 corrections, prompted by a second round of review of v1.1,
were equally substantive: replacing the $\kappa_{\mathcal X}$ third term
from set-cardinality of mapping spaces to \emph{topological weight} of the
loop groups (Convention \ref{conv:kappa} and the new explanatory paragraph),
which restores compatibility between the definition and the geometric
attestations of \S\ref{sec:direct}; correcting the bibliographic
identifier and journal reference for Campion-Cousins-Ye (the v1.1
preprint cited an unrelated arXiv identifier; v1.2 cites the published
JSL paper, arXiv:1808.04915); correcting the v1.1 claim that the
ACF$_0$ instance was ``not attained'' (it does attain the envelope at
$\kappa = \aleph_0$, since both sides equal $\continuum$); strengthening
the model-theoretic context with the classical bound $|\GalL(T)| \le 2^{|T|}$
\eqref{eq:lascar-bound} and direct citations to Campion-Cousins-Ye for
the DLO and ACF examples; and tightening the cardinality computations
in \S\ref{sec:direct} to use weight rather than underlying cardinality.

The remaining errors are our responsibility.

The axiom-explorer harness is open source; reproducibility instructions and
full per-phase reports are at \cite{AxiomExplorerRepo}.

\section{Reproducibility}

All artefacts of the experimental run that produced this preprint are
versioned at \cite{AxiomExplorerRepo}:

\begin{itemize}
\item Full per-phase reports under \texttt{research/phase\{0..5\}-*}.
\item Python harness, finite-model sanity checks under
\texttt{src/axiom\_explorer/}.
\item Lean 4 axiomatic skeleton at
\texttt{lean/AxiomExplorer/SyntheticStoneDuality.lean}, builds clean with
\texttt{leanprover/lean4:v4.16.0}.
\item Bibliometric raw data under \texttt{data/phase0/raw/}.
\item Cited papers can be re-downloaded by \texttt{curl} from arXiv; paths
are documented in the repository's \texttt{HANDOFF.md}.
\item Test suite (51 tests, all green at the time of writing) is the
reproducibility floor for the finite computational claims of this preprint.
\end{itemize}

\section{Invitation}

This v1.1 is intended as a \emph{registered conjectural observation} and an
explicit invitation to specialists in any of the involved branches. The four
responses we are most interested in, in decreasing order of preference:

\begin{enumerate}
\item ``This formulation is folklore in [community~$X$], see [reference~$Y$].''
\item ``There is a counterexample to the upper bound,
  namely~$[Z]$.''
\item ``The upper bound is plausible but the saturation question requires the
  following sharpening: \ldots''
\item ``I had not seen the conjecture stated this way; here is a specific
  place where it would help / fail.''
\end{enumerate}

Any of the four is a useful update. We will issue a v1.2 of this preprint if
substantive feedback is received.

\appendix

\section{Higher \texorpdfstring{$\pi_n$}{pi-n} extension and the K3 question}\label{app:higher}

\emph{This appendix is L3 (speculative). It is appendix material, not body
content, in response to a reviewer's observation that the higher-$\pi_n$
extension was insufficiently supported in the v1 main body.}

A naive structural argument suggests that each
$\pi_n(\Bcat\,\Pt^{\mathrm{coh}}(\mathcal X))(\ast)$ is a subquotient of
profinite groups indexed by points of $\mathcal X$, and would therefore obey
$\le 2^{\kappa_{\mathcal X}}$ for all $n \ge 1$. We do \emph{not} have a proof
or even a reference for this naive argument, only the analogy with $n = 1$.
Pending such a proof, we tentatively pose:

\begin{question}[Higher envelope, conditional]\label{q:higher}
For $\mathcal X$ spectral and $n \ge 1$, does
$| \pi_n(\Bcat\,\Pt^{\mathrm{coh}}(\mathcal X))(\ast) | \le 2^{\kappa_{\mathcal X}}$
hold?
\end{question}

The cleanest test case we identify is the K3 surface over $\C$ at $n=2$,
which is motivated by the rank-$22$ second integral cohomology
$H^2(\mathrm{K3}, \mathbf{Z}) \cong \mathbf{Z}^{22}$ (the K3 lattice, of
signature $(3,19)$) and the corresponding rank-$22$ $\ell$-adic \'etale
realisation. The natural question, set against the conjectural higher
envelope above, is:

\begin{question}[K3 at $n=2$]
Does $|\pi_2^{\mathrm{cond}}(\mathrm{K3}_{\C}, \overline{x})(\ast)|$ attain
$2^{\continuum}$, or does it remain bounded by the rank-$22$ classical
structure?
\end{question}

We do not know of a published computation of $\pi_n^{\mathrm{cond}}$ for K3
surfaces at any $n \ge 2$. The question is recorded here as a forward target
for future iterations or expert input.

\bibliographystyle{plain}
\bibliography{references}

\end{document}
